\DocumentMetadata{
  pdfversion=2.0,
  pdfstandard=ua-2,
  lang=en-US,
  tagging=on,
  tagging-setup={math/setup=mathml-SE,tabsorder=structure}
}
\documentclass[11pt,letterpaper]{article}
\usepackage[margin=0.68in,includefoot]{geometry}
\usepackage{newcomputermodern}
\usepackage{amsmath,amscd,mathtools}
\usepackage{tikz,adjustbox}
\providecommand{\square}{\mdlgwhtsquare}
\usepackage[hidelinks]{hyperref}
\hypersetup{
  pdftitle={Complex Analysis PhD Exam, September 2009},
  pdfauthor={University of Florida Department of Mathematics},
  pdfsubject={Graduate mathematics examination},
  pdfdisplaydoctitle=true,
  bookmarksopen=true
}
\setcounter{secnumdepth}{0}
\setlength{\parindent}{0pt}
\setlength{\parskip}{0.28em}
\setlength{\emergencystretch}{3em}
\setlength{\leftmargini}{2.15em}
\setlength{\leftmarginii}{2.65em}
\setlength{\leftmarginiii}{3em}
\pagestyle{empty}
\raggedbottom
\allowdisplaybreaks
\definecolor{ProblemConcern}{HTML}{7A1F1F}
\newenvironment{problemconcern}{%
  \par\tagstructbegin{tag=Aside}\begingroup\color{ProblemConcern}
}{%
  \par\endgroup\tagstructend\nopagebreak[4]\vspace{0.12em}
}
\NewTaggingSocketPlug{block/list/label}{examproblem}{%
  \tagstructbegin{tag=Lbl}\tagstructend%
  \tagstructbegin{tag=\LBody}%
  \tagstructbegin{tag=\ProblemHeadingTag,title-o={\ProblemHeadingText}}%
  \tagmcbegin{tag=\ProblemHeadingTag,actualtext={\ProblemHeadingText}}%
  #2%
  \tagmcend%
  \tagstructend%
}
\newenvironment{examproblems}[2]{%
  \def\ProblemHeadingTag{#2}%
  \AssignTaggingSocketPlug{block/list/label}{examproblem}%
  \begin{enumerate}%
  \setcounter{enumi}{#1}%
  \renewcommand{\labelenumi}{\textbf{\theenumi.}}%
  \setlength{\topsep}{0.35em}%
  \setlength{\partopsep}{0pt}%
  \setlength{\itemsep}{0.48em}%
  \setlength{\parsep}{0.18em}%
}{\end{enumerate}}
\newenvironment{examsubparts}{%
  \AssignTaggingSocketPlug{block/list/label}{default}%
  \begin{enumerate}%
  \setlength{\topsep}{0.18em}%
  \setlength{\partopsep}{0pt}%
  \setlength{\itemsep}{0.16em}%
  \setlength{\parsep}{0.1em}%
}{\end{enumerate}}
\newenvironment{examinstructions}{%
  \par\begingroup\itshape
}{%
  \par\endgroup\vspace{0.18em}
}
\makeatletter
\renewcommand\subsection{\@startsection{subsection}{2}{0pt}%
  {-1.25ex plus -.25ex minus -.1ex}%
  {0.55ex plus .1ex}%
  {\normalfont\large\bfseries}}
\renewcommand{\@maketitle}{%
  \newpage\null\vskip -1em%
  {\footnotesize\bfseries
    \href{https://www.ufl.edu/}{University of Florida}, \href{https://math.ufl.edu/}{Department of Mathematics}\par}%
  \vskip -0.5em%
  \tagstructbegin{tag=H1,title={Complex Analysis PhD Exam, September 2009}}%
  \tagpdfparaOff%
  \tagmcbegin{tag=H1}%
    {\LARGE\bfseries\href{https://gma.math.ufl.edu/exams/complex-analysis-phd/}{Complex Analysis PhD Exam}, September 2009\par}%
  \tagmcend%
  \tagpdfparaOn%
  \tagstructend%
  \vskip 0.25em%
  \tagmcbegin{artifact}%
    \hrule height 1pt%
  \tagmcend%
  \vskip 1em%
}
\makeatother
\title{Complex Analysis PhD Exam, September 2009}
\author{}
\date{}
\begin{document}
\maketitle
\thispagestyle{empty}
\begin{examproblems}{0}{H2}
\def\ProblemHeadingText{Problem 1.}
\item
A fixed point of a function~\(f\) is a point~\(x\) with \(f(x)=x\). Let~\(G\) be an open, simply connected subset of \(\mathbb C\). If~\(f\) is analytic on~\(G\), \(f(G)\subset G\), and~\(f\) has two fixed points, show that~\(f\) is the identity map.
\def\ProblemHeadingText{Problem 2.}
\item
Let~\(G\) be an open, connected subset of \(\mathbb C\) and~\(f\) be analytic on~\(G\). Note that we are not assuming that~\(G\) is simply connected.
\begin{examsubparts}
\item[\textnormal{(a)}]
Assume now that for another open, connected set~\(H\), \(f(G)\subset H\) and \(h:H\to\mathbb R\) is harmonic. Show that \(h\circ f\) is harmonic.
\item[\textnormal{(b)}]
Prove or disprove: Both the real and imaginary parts of~\(f\) are harmonic functions.
\item[\textnormal{(c)}]
Prove or disprove, if \(u:G\to\mathbb R\) is harmonic, then there is a harmonic function \(v:G\to\mathbb R\) so that \(u+iv\) is analytic.
\end{examsubparts}
\def\ProblemHeadingText{Problem 3.}
\item
Give an explicit biholomorphism from the “wedge” \(\{z:0<\arg(z)<\pi/2\}\) onto the open unit disk.
\def\ProblemHeadingText{Problem 4.}
\item
Evaluate the integral, and justify each step.
\[
\int_{-\infty}^{\infty}\frac{x\sin x}{(x^2+1)^2}\,dx.
\]
\def\ProblemHeadingText{Problem 5.}
\item
Let~\(f\) be an entire function of finite order. Prove that if the order is not an integer, then~\(f\) must have infinitely many zeros. Does there exist an entire function of infinite order with finitely many zeros? Explain.
\def\ProblemHeadingText{Problem 6.}
\item
\begin{problemconcern}
\textbf{Warning: Suspected error.} The statement is preserved as printed, but it overlooks the possible case \(g\equiv0\), for example when \(g_n\equiv1/n\). The intended correction is not uniquely determined: one could assume \(g\not\equiv0\) or conclude that~\(g\) is identically zero or has only real zeros.
\end{problemconcern}
Assume each \(g_n\) is entire and \(g_n\to g\) uniformly on compact sets in \(\mathbb C\) (i.e. \(g_n\to g\) in \(H(\mathbb C)\)). If each \(g_n\) has only real zeros, show that~\(g\) has only real zeroes.
\def\ProblemHeadingText{Problem 7.}
\item
Assume that~\(f\) and~\(g\) are analytic and nonvanishing on \(\{z\in\mathbb C:|z|<2\}\) and that \(|f(z)|=|g(z)|\) when \(|z|=1\). Show that there is a constant \(\lambda\in\mathbb C\) with \(|\lambda|=1\) and \(f(z)=\lambda g(z)\) for all \(z\in B_2(0)\).
\def\ProblemHeadingText{Problem 8.}
\item
Assume that~\(f\) is entire, and for some integer \(n>0\),
\[
\lim_{z\to\infty}\frac{f(z)}{z^n}=c,
\]
 for some finite, nonzero complex number~\(c\). What can you say about~\(f\)?
\end{examproblems}
\end{document}
